\documentclass[11pt]{article}
\usepackage[margin=1.1in]{geometry}
\usepackage{amsmath,amssymb}
\usepackage{booktabs}
\usepackage{array}
\usepackage[hidelinks]{hyperref}
\usepackage{xcolor}
\usepackage{parskip}

\newcommand{\CNOT}{\textsc{cnot}}
\newcommand{\CCNOT}{\textsc{ccnot}}
\newcommand{\code}[1]{\texttt{#1}}
\newcommand{\Sone}{S_1}

\title{\textbf{New Gadgetize}\\[2pt]
\large The $\times 2$ gadgetization recipe: design rationale and changes}
\author{local\_mixing project}
\date{July 20, 2026 \\ \small branch \code{ssg-gen-mix-clean}, commit \code{49c8cc02}}

\begin{document}
\maketitle

\begin{abstract}
\noindent\emph{New gadgetize} is the name of the current production recipe for the
$\times 2$ (two-share) gadgetization path of the CNOT-native backend
(\code{sss --cnot --gadgetize --slice-zero-ccnot}). Relative to the previous
plain gadgetization it adds two things: a \emph{zero-slice preblock} $\Sone$
that makes the circuit compute the target function only on one designated
input slice (and, as a corollary, stops the inverse circuit from handing out
$C^{-1}$), and a leaner, faster \emph{randomization-gadget policy} that
replaces the three-move RG set by the orthogonal basis $\{RG2, RG3\}$ at four
times the rate. This note records the construction, the rationale, the
guarantees (exact and measured-heuristic), the design alternatives that were
considered and rejected, and the precise differences from previous versions.
\end{abstract}

\section{Background: the $\times 2$ gadget path before this change}

The gadget path transforms a reversible circuit $C$ on $n$ wires into a
circuit on $2n$ wires: the first half ($x$, wires $0..n{-}1$) carries data,
the second half (wires $n..2n{-}1$) is auxiliary. The emitted structure is
\[
[\text{bookend}]\;[\text{$W_i$ encode}]\;[\text{SG/RG body}]\;[\text{decode}]\;[\text{bookend}],
\]
where the bookends are $\max(2n\ln n, 64)$ random g57 gates junk-targeting
the second half; the $W_i$ sweep (7 \CNOT{}s per value) places every logical
value $v$ into a two-share pair $(\mathit{share},\mathit{pad})$ with
$v = \mathit{share}\oplus\mathit{pad}$, using the aux contents as masks; each
logical g57 gate is realized by a four-fragment masking-safe SG (a four-cube
ESOP of $B \lor \lnot C$ over the four carriers, all fragments targeting one
carrier of the output value); and the decode folds each value back onto its
wire. The output contract (``GadgetLow'') is:
\[
A(x, a) = \bigl(C(x),\; \text{junk}\bigr) \qquad \text{for \emph{every} aux value } a .
\]

Two properties of this contract motivated the change:
\begin{enumerate}
\item \textbf{No slice conditioning.} The circuit computes $C$ regardless of
the aux input; there is no notion of a distinguished slice on which the
functionality lives, unlike the feistel path's slice modes.
\item \textbf{The inverse is free.} Because the contract holds for every
$a$, the \emph{inverse} circuit satisfies $G^{-1}(p, v) = (C^{-1}(p), \cdot)$
for every junk value $v$: anyone holding the (public, mixed) circuit can
run it backwards and read $C^{-1}$ off the low wires. The forward oracle
gives $C$; the naked gadget also gives away $C^{-1}$.
\end{enumerate}

Additionally, the body's randomization used three moves (RG1/RG2/RG3,
uniform) at a rate of one RG per two SGs; \S\ref{sec:rg} argues this set was
redundant and the rate low.

\section{The zero-slice preblock $\Sone$}

\subsection{Construction}

$\Sone$ is prepended to the gadget (gate-list concatenation). It consists of
$G$ gates (default $10n$, flag \code{--slice-zero-ccnot-gates}) in
\emph{one uniform random order}, drawn from exactly two shapes, all with
positive-polarity controls and all targets on the first half:
\[
\underbrace{x_i \mathrel{\hat{}}= a_j}_{\text{\CNOT, } \approx 1/3}
\qquad\qquad
\underbrace{x_i \mathrel{\hat{}}= x_j \wedge a_k}_{\text{\CCNOT, } \approx 2/3}
\]
(one control of each \CCNOT{} on each half). Every gate reads at least one
second-half wire, so on the all-zero slice $a=0$ every gate is
\emph{individually} dead: $\Sone(x,0)=(x,0)$ gate by gate, with no ordering
or pairing argument. The second half passes through untouched (targets are
confined to the first half), so the slice survives on its wires into the
gadget.

For a fixed slice $a\neq 0$, each \CCNOT{} collapses to a within-$x$ \CNOT{}
(when its aux control is hot) and each \CNOT{} to a constant flip, so the
disturbance is a slice-indexed invertible \emph{affine} map
\[
M_a(x) \;=\; L_a\,x \;\oplus\; t_a .
\]
The composite circuit therefore computes
\[
A(x, 0) = \bigl(C(x), J\bigr), \qquad
A(x, a) = \bigl(C(M_a(x)), J'\bigr) \;\; (a \neq 0),
\]
and since $C$ is a permutation, a wrong slice reproduces $C$ at $x$ exactly
when $M_a$ fixes $x$.

\subsection{Guarantees}

\emph{Superseded 2026-07-25; see} \texttt{docs/BAND\_HARDENING} \S3. The
construction below is the original one, whose guarantee was exact on slices
firing no \CCNOT{} and heuristic elsewhere. It has since been rebuilt around a
target/control split whose order-free algebra makes ``only the all-zero slice
is fixed'' a \textbf{theorem for every slice}, checked outright on the emitted
gates by a rank test---and the slice was widened to include the product-share
band. The historical statement follows.

\textbf{Exact (original).} The \CNOT{} pin multiset (a random permutation seed
plus random extras) was resampled until its target-by-control parity matrix
$\mathbf{C}$ was invertible over $\mathrm{GF}(2)$. Consequently, on any slice
that fires no \CCNOT{} the disturbance is the pure translation
$x \oplus \mathbf{C}a \neq x$: every such slice is disturbed at every input,
provably. (Invertibility also guarantees every aux wire is read.)

\textbf{Heuristic, measured (original).} On slices firing \CCNOT{}s, ``no wrong
slice is fixed'' was not a theorem: it fails if the fired transvection
subsequence composes to the identity \emph{and} the interleaving-conjugated
\CNOT{} translations cancel. This was measured directly (exhaustive slice check
per sampled preblock, gate budget $10n$):

\begin{center}
\begin{tabular}{lccccc}
\toprule
$n$ & 3 & 4 & 5 & 6 & 8 \\
\midrule
preblocks with a fixed wrong slice &
$\approx 4\cdot 10^{-3}$ & $4\cdot 10^{-4}$ & $2\cdot 10^{-5}$ & $4\cdot 10^{-5}$ & $0$ \\
trials & $2{,}000$ & $50{,}000$ & $50{,}000$ & $50{,}000$ & $50{,}000$ \\
\bottomrule
\end{tabular}
\end{center}

The rate decayed rapidly in $n$ and was negligible at production widths. The
current construction removes the residual entirely: exactness now holds at
every $n$, including the toy widths where this table shows failures, and the
test suite asserts it exhaustively rather than pinning a counterexample.

\textbf{Inverse protection.} In the inverse circuit $A^{-1}$ (the reversed
gate list; every gate is an involution), $\Sone^{-1}$ runs \emph{last} and
fires on the gadget's generically nonzero mask residue, junking the low
half. This closes the free-$C^{-1}$ hole of \S1: a bare gadget's inverse
returns $C^{-1}$ for any junk input; with $\Sone$ it does not. This is
regression-tested
(\code{slice\_block\_stops\_the\_inverse\_from\_revealing\_c\_inverse}).

\subsection{Why this vocabulary}

The preblock is pure positive-polarity \CNOT{}/Toffoli material:
indistinguishable in gate shapes from ordinary mixed-circuit content, with
no complemented conjunctions and no polarity pattern encoding a slice. Two
contrasts drove this choice:

\begin{itemize}
\item The feistel path's random-slice mode (mode~3) encodes its public slice
\emph{in the clear} in the control polarities of its core OR-gates. With the
slice fixed to the all-zero vector and positive controls throughout, there
is nothing to read off.
\item The legacy r57-only vocabulary ($a \mathrel{\hat{}}= b \lor \lnot c$)
cannot express a gate that both reads an $x$ wire and is dead on a slice;
its slice constructions needed commuting-pair cancellation tricks. The
\CCNOT{} shape $x_i \mathrel{\hat{}}= x_j \wedge a_k$ achieves per-gate
deadness and $x$-fanout simultaneously.
\end{itemize}

\subsection{Alternatives considered and rejected}
\label{sec:rejected}

\textbf{Contiguous \CNOT{} certificate block (rigorous uniqueness).} Placing
all \CNOT{}s in one contiguous run at an end of the block makes uniqueness a
theorem: the run commutes into $x \oplus \mathbf{C}a$, the \CCNOT{} soup is
translation-free per slice, and $M_a = \mathrm{id}$ forces
$\mathbf{C}a = 0$, i.e.\ $a=0$. This was implemented (in one-end and
symmetric two-end variants) and \emph{rejected by design decision}: the
uniform random order was preferred --- the \CCNOT{}s do not commute, so the
order itself is part of the randomness, and bunched \CNOT{} runs are
structure the artifact should not carry. The cost of that choice is exactly
the heuristic table above.

\textbf{Two-slice sandwich $\Sone \,\|\, G \,\|\, S_2$.} A closing slice
block was fully built (with the gadget re-plumbed to decode onto the high
wires so $S_2$'s first-half targets would only touch junk) and then
\emph{reverted}: it belongs to the feistel-style sandwich construction
(where an $N$-step moves the answer to the second half and a random $D$
junks the first), not to the general $\times 2$ gadget mechanism, whose
answer exits on the low wires --- there, a closing $S_2$ ruins correctness.
What was kept from the detour is the inverse-direction analysis and its
test.

\section{The new RG policy}
\label{sec:rg}

\subsection{The moves}

With two-share pairs $i=(i_0,i_1)$ and $j=(j_0,j_1)$:

\begin{center}
\begin{tabular}{llll}
\toprule
 & gates & effect & role \\
\midrule
RG1 & 6 \CNOT{}s & values swap wire-sets & \emph{dropped} \\
RG2 & 3 \CNOT{}s (swap $i_0 \leftrightarrow j_0$, re-pair crosswise) & pairing re-dealt, values intact & sole re-pairer \\
RG3 & 2 \CNOT{}s ($i_0 \mathrel{\hat{}}= r$, $i_1 \mathrel{\hat{}}= r$, foreign $r$) & masks refreshed, value intact & sole mask injector \\
\bottomrule
\end{tabular}
\end{center}

All moves are linear and \emph{gate-locally non-complete}: no physical gate
ever reads both carriers of one logical value, so first-order prefix
masking is preserved --- this invariant is unchanged from the previous
version and regression-tested.

\subsection{Why $\{RG2, RG3\}$ is a basis and RG1 is redundant}

RG2 and RG3 are orthogonal: RG2 changes \emph{which wires are paired}
(compositions of RG2 only \emph{permute} wire contents --- each application
is a content swap), while RG3 is the only move that \emph{recombines mask
values} (XORing a foreign carrier into both members of a pair). Neither can
do the other's job, and together they generate every sharing configuration.

RG1 is a composite: two generalized RG2s (swap $i_0\leftrightarrow j_0$,
then $i_1\leftrightarrow j_1$) reproduce its value-swap at identical
\CNOT{} cost, and the one thing the swap route lacks --- RG1's mask
recombination (its outputs are XOR mixtures such as
$(b\oplus c,\, b\oplus d,\, a,\, b)$, versus the pure content permutation
$(d,c,b,a)$ of the two-swap route) --- is precisely RG3's job. RG1 earned
its former place by packaging value-swap-plus-mask-stir in 6 \CNOT{}s; with
RG3 present at a higher rate it is redundant, and dropping it simplifies the
move algebra.

\subsection{Rate}

Previously: one uniform RG1/RG2/RG3 draw after every second SG
($0.5$ RGs per SG, $\approx 1.8$ \CNOT{}s per SG on average). Now:
\code{--rg-frequency} independent uniform RG2/RG3 draws \emph{between every
two consecutive SGs} (none after the last, immediately before the decode).
The default of 2 gives 2 RGs per gap $= 4\times$ the move rate,
$\approx 5$ \CNOT{}s per SG ($2 \times \frac{3+2}{2}$) $\approx 2.7\times$
the \CNOT{} volume. The sharing now drifts faster than any single SG's
four-fragment footprint. Note the flag's meaning \emph{changed} for this
path (it was ``one RG every $K$ SGs''); the feistel and legacy paths keep
the old meaning and the full RG1/RG2/RG3 set.

One acknowledged trade-off: RG2 is a contiguous 3-\CNOT{} swap, so the raw
(unmixed) artifact contains more recognizable swap blocks than RG1's
interleaved networks produced. The mixer dissolves these immediately; the
previous version's material was prettier pre-mixing, not stronger
post-mixing.

\section{Differences from previous versions: summary}

\begin{center}
\renewcommand{\arraystretch}{1.25}
\begin{tabular}{p{3.4cm}p{4.9cm}p{5.6cm}}
\toprule
 & \textbf{plain gadgetize} (previous) & \textbf{new gadgetize} (\code{49c8cc02}) \\
\midrule
invocation & \code{sss --cnot --gadgetize} & \code{sss --cnot --gadgetize --slice-zero-ccnot} \\
contract & $C(x)$ on low wires for \emph{every} aux input & $C(x)$ on low wires \emph{only} on the $a{=}0$ slice; $C(M_a(x))$ off-slice \\
inverse circuit & returns $C^{-1}$ for any junk input & low half junked by $\Sone^{-1}$ \\
preblock & none & $10n$-gate \CNOT/\CCNOT{} slice block, uniform order, invertible parity matrix \\
RG set & RG1, RG2, RG3 (uniform) & RG2, RG3 only \\
RG rate & one per 2 SGs ($\approx 1.8$ \CNOT/SG) & two per SG gap ($\approx 5$ \CNOT/SG) \\
\code{--rg-frequency} & SGs per RG & RGs per SG gap (this path only) \\
verification & GadgetLow over random aux & GadgetLow pinned to the zero slice, at transform and every round \\
size ($n{=}8$, 40 gates) & $\approx 560$ gates & $\approx 670$ gates \\
\bottomrule
\end{tabular}
\end{center}

Relation to the feistel slice modes (for orientation): the feistel path
offers a linear zero-slice mode ($x \oplus A(y{\lor}z)$), a nonlinear
hardcoded zero-slice mode, and a random \emph{public} slice mode
($R^{-1}BR$, uniquely fixing by construction but carrying the slice in its
gate polarities). New gadgetize is the $\times 2$ analogue of a zero-slice
mode with mixing-native vocabulary: per-gate deadness, no polarity
signature, exact uniqueness on the \CCNOT{}-silent subclass and measured
heuristic uniqueness elsewhere.

\section{Open questions}

\begin{enumerate}
\item \textbf{$\Sone$ randomizer strength.} Per slice, $\Sone$ is a shallow
affine map at the $10n$ budget (each first-half wire written $\approx 10$
times; no nonlinearity in $x$, since $x{\wedge}x$ controls are excluded by
the gate-shape spec). Whether this matters for post-mixing inversion
resistance is open; the knob is \code{--slice-zero-ccnot-gates}, and a
planned measurement is avalanche/dislocation of an isolated $\Sone$ versus
budget at production $n$.
\item \textbf{Feistel port.} Neither the RG policy nor the slice block has
been ported to the feistel path. Note the invertibility certificate does
not transfer as-is: the feistel aux block has $2n$ wires, so the parity
matrix would be $n \times 2n$ and never injective; a feistel variant needs
mode-3-style pairing or a purely heuristic argument.
\item \textbf{Deployment.} Not yet built on the .242 server.
\end{enumerate}

\vspace{6pt}
\noindent\textbf{Provenance.} Code: \code{slice\_zero\_ccnot\_preblock},
\code{gadgetize\_with\_slice\_zero\_ccnot}, and the RG loop in
\code{src/replace/gadgets.rs}; driver in
\code{src/replace/main\_mix\_cnot.rs}. Tests: four slice-block unit tests
plus the zero-round driver case; full suite 139 passing at commit
\code{49c8cc02}. Companion quick-reference: \code{docs/NEW\_GADGETIZE.md}.

\end{document}
